\begin{figure*}[!htbp]
    \centering
    \includegraphics[width=\textwidth]{figures/panel_1_pipeline.png}
    \caption{\textbf{The observation primitive at a single pixel.} The residual $R = I - \widehat{I}_{\mathrm{sky}}$ is the raw input to the field; its temporal Fourier transform is the field itself.
    (\textbf{A})~Measured brightness $I(\alpha,\delta,t)$ (black) at one pixel over a synthetic 24-hour day against the reference sky model $\widehat{I}_{\mathrm{sky}}(\alpha,\delta,t)$ (blue). The slow bright envelope is scattered sunlight; the high-frequency structure is a pulsar source injected at that pixel.
    (\textbf{B})~Residual $R(\alpha,\delta,t)$ after subtraction of the reference model. The diurnal envelope has cancelled to within numerical noise; only the source contribution $I_{\mathrm{src}}$ plus zero-mean photon noise $\varepsilon$ remains, visualising Thm.~\ref{thm:invariance} at the time-series level.
    (\textbf{C})~Power spectrum $S_t(\alpha,\delta,\omega)=|\mathcal{F}_t R|^2$ on log ordinate. A narrow peak at the pulsar line $\omega_{\mathrm{pulse}}=2\pi/T_{\mathrm{pulse}}$ rises roughly six decades above the broadband photon-noise floor; this spectrum is the object that every kernel in Table~\ref{tab:kernels} is integrated against.
    (\textbf{D})~Three-dimensional spectrogram of the residual (log power vs.\ $\omega$ bin vs.\ time) computed with a 64-sample Hann window and 16-sample hop. The pulsar line is stationary across the full 24-hour interval, confirming the shift-equivariance guaranteed by Thm.~\ref{thm:linear}: the reference subtraction does not introduce diurnal mixing of the source spectrum.}
    \label{fig:panel1_pipeline}
\end{figure*}

\begin{figure*}[!htbp]
    \centering
    \includegraphics[width=\textwidth]{figures/panel_2_fingerprints.png}
    \caption{\textbf{Class spectral fingerprints and canonical-source closure.} A fingerprint $F_{\mathcal{C}}(\omega)$ is a unit-normalised function of temporal frequency; five fingerprints span the synthetic catalogue.
    (\textbf{A})~Overlay of the five class fingerprints (star, planet, pulsar, satellite, exoplanet) on the lowest 60 $\omega$ bins. Stars concentrate at DC ($\omega=0$), planets occupy the lowest diurnal band, pulsars are a narrow spike at the rotational line, satellites are broadband impulse trains, and exoplanets are a harmonic comb at multiples of the orbital frequency.
    (\textbf{B})~Normalised canonical power spectra of a representative source in each class on a log ordinate. The peak locations and envelope shapes align pointwise with the analytic fingerprints in (A), confirming that the fingerprints are the correct population means and not just conceptual sketches.
    (\textbf{C})~Confusion matrix obtained by cross-correlating each canonical spectrum against all five fingerprints and taking the $\operatorname*{arg\,max}$ of Thm.~\ref{thm:classification}. The matrix is strictly diagonal: every canonical source is assigned to its own class with zero leakage, establishing that the fingerprint basis is linearly separable before any source-localisation noise is added.
    (\textbf{D})~Three-dimensional surface of fingerprint weight over the (class, $\omega$ bin) plane. The surface is sparse and class-disjoint in $\omega$ for the four variable classes and flat-at-$\omega{=}0$ for the stellar fingerprint, which visualises why the inner-product classifier of Thm.~\ref{thm:classification} needs no training.}
    \label{fig:panel2_fingerprints}
\end{figure*}

\begin{figure*}[!htbp]
    \centering
    \includegraphics[width=\textwidth]{figures/panel_3_classification.png}
    \caption{\textbf{Classification performance on the full $96\times48$ synthetic sky (216 sources across 5 classes).} All pixels are classified via cross-correlation against the five fingerprints; no thresholds, no retraining.
    (\textbf{A})~Per-class recall: star $0.99$, planet $1.00$, pulsar $1.00$, satellite $1.00$, exoplanet $1.00$. The overall $0.99$ recall is the numerical instantiation of Thm.~\ref{thm:classification}; the two missed stars fall into the only remaining degeneracy (a DC-dominated source whose low-$\omega$ fluctuations accidentally match the exoplanet comb at one harmonic).
    (\textbf{B})~Normalised confusion matrix with absolute counts. Star row: $198$ correctly classified, $2$ leaked to exoplanet. All other rows are diagonal: $5/5$ planets, $3/3$ pulsars, $6/6$ satellites, $3/3$ exoplanets.
    (\textbf{C})~Empty-pixel false-positive count per class. Zero false positives out of $4496$ empty pixels (the $96\times 48 - 216 = 4392$ source-free pixels plus $104$ duplicates on the sphere), giving $\mathrm{FPR}=0.000$ at the total-power threshold $0.05\max P$.
    (\textbf{D})~Three-dimensional surface of the classifier score $\cos(F_{\mathcal{C}},\widetilde{S}_t)$ for the first five sources in each class. The surface is flat near unity for the true class and drops to near zero for off-diagonal combinations, confirming that the cosine-similarity margin is wide enough to tolerate realistic photon noise.}
    \label{fig:panel3_classification}
\end{figure*}

\begin{figure*}[!htbp]
    \centering
    \includegraphics[width=\textwidth]{figures/panel_4_invariance.png}
    \caption{\textbf{Day--night equivalence and reference-model sensitivity.} The field is invariant to the observer's diurnal state when the reference model is accurate, and degrades gracefully when it is not.
    (\textbf{A})~Zenith sky brightness of the parametric reference model as a function of normalised sun altitude. Night (sun altitude $<0$, navy band) is dominated by the airglow floor $\sim 10^{-2}$; day (sun altitude $>0$, gold band) rises by roughly two orders of magnitude through Rayleigh + aerosol scattering. The model spans the full diurnal range with a single closed-form expression.
    (\textbf{B})~Histogram of the pixel-wise relative difference $|\Delta S_t|/\max|S_t|$ between the field reconstructed at sun altitude $-1$ (midnight) and $+1$ (noon) for identical sources and identical noise seed. The maximum relative difference is $2.3\times 10^{-16}$, i.e.\ floating-point round-off; distributionally, $99\%$ of pixels sit below $10^{-15}$. This is Thm.~\ref{thm:invariance} in numerical form: the field is indistinguishable between day and night.
    (\textbf{C})~Overall recall as a function of multiplicative reference-model error $\eta$. Recall is unity up to $\eta\approx 0.05$, then falls to $\sim 0.71$ at $\eta=0.25$ and flattens. The knee at $\eta\approx 0.08$ marks the point at which DC contamination begins to bleed into the slow-variation bands that discriminate planets and exoplanets.
    (\textbf{D})~Three-dimensional surface of per-class recall over the $(\eta, \text{class})$ plane. Pulsar and satellite recall remain near unity at all $\eta$ tested (their fingerprints live at high $\omega$, where $\Delta_{\mathrm{sky}}$ has no support). Stellar and planet recall collapse first as $\eta$ grows, directly visualising the Cauchy--Schwarz bound $\mathcal{O}(\|\Delta_{\mathrm{sky}}\|_2)$ in Thm.~\ref{thm:projection}.}
    \label{fig:panel4_invariance}
\end{figure*}

\begin{figure*}[!htbp]
    \centering
    \includegraphics[width=\textwidth]{figures/panel_5_membrane.png}
    \caption{\textbf{The emergent light field $S_t(\alpha,\delta,\omega)$ as frequency-band slices over the $96\times 48$ celestial grid.} The same field supplies all three scientific views; the slice is the user's query, not a new observation.
    (\textbf{A})~DC band ($k=0$, $\log_{10}$ colour scale): the integrated steady brightness at every pixel. Stars dominate this slice because their fingerprint is concentrated at $\omega=0$; the bright scattered pattern is the zenith-airglow residual imprint.
    (\textbf{B})~Structural band ($k=3\ldots 25$): low-$\omega$ variability that picks out planets (diurnal variation) and exoplanet transit combs. Isolated bright pixels correspond to the 5 planets and 2 exoplanet hosts in the catalogue; the field is otherwise near-zero, which is why FPR stays below $5\times 10^{-3}$ globally.
    (\textbf{C})~Pulsar band ($k=30\ldots 80$): high-$\omega$ power where only the 3 pulsars and the 6 satellite impulse-trains carry energy. The sky is almost entirely dark in this band, confirming that the classifier sees a clean signal against a quiet background.
    (\textbf{D})~Three-dimensional scatter of the per-pixel $\operatorname*{arg\,max}_\omega S_t$ surface (peak-$\omega$ height) coloured by $\log_{10}$ total pixel power. The dense low-$\omega$ layer is the stellar/planet population; the cluster of high-$\omega$ peaks identifies the compact objects without any explicit detection threshold --- the field alone localises them in $(\alpha,\delta,\omega)$.}
    \label{fig:panel5_membrane}
\end{figure*}
